\documentclass[12pt,a4paper]{article}
\usepackage[slovene]{babel}
\usepackage[utf8]{inputenc}
\usepackage{amsmath}
\usepackage{amsthm}
\usepackage{graphicx}
\usepackage{mdframed}
\usepackage{float}
\usepackage{url}
\usepackage{authblk}

\theoremstyle{definition}
\newenvironment{zgled}{%
  \begin{mdframed}[linewidth=0.5pt, topline=true, bottomline=true, leftline=true, rightline=true, 
                  innertopmargin=\baselineskip, innerbottommargin=\baselineskip]%
  \textbf{Zgled.}%
}{%
  \end{mdframed}%
}

\begin{document}

\setlength{\parskip}{1em}
\setlength{\parindent}{0pt}

\title{Zasebne digitalne valute - primer LIBRA}
\author{Urh Videčnik, Jure Kraševec}
\date{\today}
\affil{Fakulteta za matematiko in fiziko, Univerza v Ljubljani}
\maketitle

\begin{abstract}
    Namen naše naloge je oceniti učinkovitost ameriške in slovenske javne
    uprave. Tega smo se lotili s CCR modelom DEA analize. Zaradi oblike modela
    smo se osredotočili samo na zdravstvo. V teoretičnem
    delu smo spoznali osnovno idejo DEA analize in kako se uporablja za 
    izboljšanje učinkovitosti. Izpostavili smo tudi pomankljivosti DEA. 
    V praktičnem delu smo s pomočjo CCR indeksa in Malmquistovega indeksa
    ugotovili, da v zdravstvu v Sloveniji in ZDA produktivnost v zadnjih
    25 letih raste, najbolj produktiven čas pa je bil čas epidemije COVID-19. 
\end{abstract}

\newpage

\begingroup
\setlength{\parskip}{5pt}
\tableofcontents
\endgroup

\newpage

\section{Uvod}

\end{document}